\section{Running RECLAIM}
\label{app:running}

This appendix records what the supplementary material contains and the exact
commands that serve a model, run one paper, run a tier, grade a run, and
regenerate the appendix listings.

\subsection{What the supplementary material contains}
\label{app:running-contents}

% src: README.md:38-56 (Layout table)
The harness is submitted as anonymized supplementary material. It is one
repository holding the dataset pipeline, the reproduction agent of
Appendix~\ref{app:harness}, the auditor of Appendix~\ref{app:rubric}, the
serving layer both attach to, and this paper's sources.
Table~\ref{tab:supplement} lists the top-level modules and what each owns. The
split is by responsibility, so a reader checking one stage of the pipeline
opens one directory.

% src: README.md:38-56
\begin{table}[h]
\caption{Top-level modules of the submitted harness.}
\label{tab:supplement}
\begin{center}
\small
\begin{tabular}{@{}lp{8.6cm}@{}}
\toprule
Path & What it owns \\
\midrule
\texttt{src/reprocli\_data/} & Dataset pipeline. arXiv e-print sources and
OpenReview supplements become one row per paper. \\
\texttt{src/reprocli\_vllm/} & Auditor core. Prompt and rubric assembly,
run-directory tools, tool loop, verdict schema, deterministic scoring, the
H100 band ladder, audit-pool selection. \\
\texttt{src/run\_arxiv\_prompt\_vllm.py} & Auditor entry point
(\texttt{--mode audit}) against an OpenAI-compatible endpoint. \\
\texttt{src/reprocli\_repro/} & Reproduction agent. Episode loop, Apptainer
sandbox, just-in-time GPU sessions with compute metering, evidence bundle,
report schema, telemetry. \\
\texttt{src/reprocli\_claude/} & The same auditor on the Anthropic Messages
API, with prompt caching and a schema-bound final turn. This is the pinned
grader. \\
\texttt{src/reprocli\_openai/} & Web-search re-check of the audit pool's
no-code Reimplement rows. \\
\texttt{src/reprocli\_serve/} & vLLM launcher. Per-model serve profiles,
launch flags, health wait, endpoint publication. \\
\texttt{scripts/serve/},\newline\texttt{scripts/reproduce/} & Slurm batch
jobs. Serve a model, and walk a tier's papers through the paired
reproduce-then-grade flow. \\
\texttt{scripts/cluster/} & The sandbox image build. \\
\texttt{tools/verify\_app/} & Static web app for the human verification pass
over the classifier's artifact verdicts. \\
\texttt{tools/run\_viewer/} & Static web app for run telemetry, transcripts,
and audits. \\
\texttt{tasks/} & Runbooks for serving each model and for the H100 recompute
pass. \\
\texttt{prompts/}, \texttt{rubric\_audit.md} & The frozen prompt templates and
the frozen audit rubric. \\
\texttt{tests/} & The offline test suite. \\
\bottomrule
\end{tabular}
\end{center}
\end{table}

% src: prompts/prompt_reproduce.txt, prompts/prompt_audit.txt, rubric_audit.md,
% src: paper_latex/prompts/, src/reprocli_repro/inputs.py:42-49,
% src: src/reprocli_repro/dataset.py:1-32
Four groups of artifacts back the numbers in this paper. The frozen prompt
templates and the frozen rubric sit at the repository root under
\texttt{prompts/} and \texttt{rubric\_audit.md}, and the copies rendered in
Appendix~\ref{app:prompts} and Appendix~\ref{app:rubric} are the same files
up to the disclosed substitution of the benchmark's name.
The Slurm batch scripts under \texttt{scripts/} are the exact jobs that
produced every sweep. The per-run evidence bundles of
Appendix~\ref{app:report-format} carry one directory per run with its
\texttt{report.json}, its \texttt{evidence/}, and its \texttt{stats.json}. The
dataset splits are packaged for the \texttt{datasets} library with a
\texttt{test} split of 100 papers and a \texttt{validation} split of 14, and
the harness loads them by repository id (Appendix~\ref{app:dataset-schema}). A reader can therefore check a
verdict against the run it graded without rerunning anything.

\subsection{Prerequisites}
\label{app:running-prereqs}

% src: requirements.txt, scripts/reproduce/dsv4_flash/easy_dsv4_flash.sbatch:259,
% src: scripts/serve/serve_gh200.sbatch:55, scripts/cluster/build_cuda_sandbox.sh:1-30
Five things have to be in place before any command below runs. The
orchestrator runs on Python 3.11, the version the batch scripts load, and its
dependencies install from \texttt{requirements.txt}. Apptainer must be
installed, and the sandbox image of Appendix~\ref{app:sandbox} must be built
once with \texttt{scripts/cluster/build\_cuda\_sandbox.sh}, which layers
\texttt{git}, \texttt{curl}, compilers, \texttt{cmake}, media libraries, a
Python 3.12 interpreter with its development headers, and \texttt{uv} onto a
CUDA base image. Slurm must be reachable with at least one GPU partition,
because \texttt{run\_gpu} allocates nodes from inside the episode. The agent
model must already be served at an OpenAI-compatible chat-completions URL,
since every client here is URL-only and hosts no model. The pinned grader
reads its credential from \texttt{ANTHROPIC\_API\_KEY}.

% src: src/reprocli_repro/cli_args.py:65-69, src/reprocli_repro/cli_args.py:148-153,
% src: scripts/cluster/build_cuda_sandbox.sh:26-28
Two environment variables set the defaults everything else inherits.
\texttt{REPRO\_WORK\_ROOT} moves the run bundles and outputs off the home
filesystem, and \texttt{REPRO\_APPTAINER\_SIF} names the sandbox image the
agent's shell steps run inside. Both have flag equivalents
(\texttt{--runs-dir}, \texttt{--apptainer-image}) that override them per run.
Listing~\ref{lst:running-prereqs} is the one-time setup.

% src: requirements.txt, scripts/cluster/build_cuda_sandbox.sh:22-28,
% src: src/reprocli_repro/cli_args.py:65-69, src/reprocli_claude/__main__.py:19-20
\begin{codebox}[lst:running-prereqs]{bash}{One-time setup (sanitized)}
$ python3 -m pip install -r requirements.txt

# Build the agent sandbox image once, then point the harness at it.
$ bash scripts/cluster/build_cuda_sandbox.sh <base-sif> <sandbox-sif>
$ export REPRO_APPTAINER_SIF=<sandbox-sif>

# Run bundles and outputs land here instead of $HOME.
$ export REPRO_WORK_ROOT=<work-root>

# Only the pinned grader needs this.
$ export ANTHROPIC_API_KEY=<key>
\end{codebox}

% src: scripts/cluster/build_cuda_sandbox.sh:110-126
The build script ends by compiling and importing a C extension inside the
finished image, so a missing header is caught at build time rather than in the
middle of a run.

\subsection{Serving an agent model}
\label{app:running-serve}

% src: src/reprocli_serve/args.py:20-60, src/reprocli_serve/launch.py:18-60,
% src: src/reprocli_serve/config.py:12-36, README.md:60-72
The serving half and the agent half talk only over a URL.
\texttt{reprocli\_serve} boots vLLM on a GPU node, binds \texttt{0.0.0.0}, and
publishes the address other nodes reach it at. Per-model launch flags live in
\texttt{src/reprocli\_serve/profiles.py}, which is the single source of truth
for tensor-parallel size, context length, KV cache dtype, the tool and
reasoning parsers, and the compilation config. Any flag passed on the command
line overrides the profile, and anything after a bare \texttt{--} is forwarded
to \texttt{vllm serve} after a compatibility filter drops flags the installed
version does not accept. Appendix~\ref{app:models} records the profile
values each evaluated model was served with.

% src: src/reprocli_serve/endpoint.py:1-45, src/reprocli_serve/config.py:12-36,
% src: src/reprocli_serve/args.py:33-56
Discovery is one small JSON file on a shared filesystem. The launcher waits
for the server's health endpoint to go green, reads a routable fabric address
off the interface named by \texttt{--iface}, and atomically writes the payload
of Listing~\ref{lst:running-endpoint} to the path given by
\texttt{--endpoint-file}. Clients resolve their base URL in a fixed priority
order: the explicit \texttt{--vllm-server-url} flag, then
\texttt{\$REPROCLI\_SERVER\_URL}, then the \texttt{base\_url} field of the
file named by \texttt{\$REPROCLI\_ENDPOINT\_FILE}. Swapping the model behind
the benchmark is therefore a change of one URL.

% src: scripts/serve/serve_gh200.sbatch:78-86, src/reprocli_serve/args.py:20-60
\begin{codebox}[lst:running-serve]{bash}{Serving one model (sanitized)}
$ PYTHONPATH=src python3 -m reprocli_serve \
    --model <hf-id-or-path> \
    --served-model-name <model-id> \
    --port 8000 \
    --tensor-parallel-size <tp> \
    --endpoint-file <endpoint-file>

# Every client on any node then finds it through this one variable.
$ export REPROCLI_ENDPOINT_FILE=<endpoint-file>
\end{codebox}

% src: src/reprocli_serve/endpoint.py:22-40
\begin{schemabox}[lst:running-endpoint]{The published endpoint file (sanitized)}
{
  "base_url": "<endpoint-url>",
  "host": "<host>",
  "port": 8000,
  "served_model_name": "<model-id>",
  "model_path": "<hf-id-or-path>",
  "node": "<node>",
  "job_id": "<job-id>",
  "health": "ready"
}
\end{schemabox}

% src: src/reprocli_serve/endpoint.py:22-40, scripts/reproduce/dsv4_flash/easy_dsv4_flash.sbatch:316-328
The \texttt{job\_id} field is what lets a resumed sweep refuse an endpoint a
dead predecessor published, so papers are not spent against an address nothing
answers at.

\subsection{Running one paper}
\label{app:running-one}

% src: src/reprocli_repro/cli_args.py:88-127, src/reprocli_repro/cli_args.py:155-175,
% src: src/reprocli_repro/dataset.py:28-35
One invocation runs one paper with one model.
\texttt{--paper-id} selects the row, and \texttt{--split} selects the
published split it is read from. The split defaults to \texttt{test}, the
frozen 100-paper benchmark, and the aliases \texttt{eval} and \texttt{dev} map
onto \texttt{test} and \texttt{validation}. \texttt{--lockfile} points at the
dataset the row comes from and accepts a Hugging Face repository id, an
\texttt{hf://} file reference, or a local JSONL path.
\texttt{--run-id} pins the run identifier, which is what makes the bundle
findable afterwards. Omitting it mints a fresh time-stamped identifier, so
repeat attempts never collide.

% src: src/reprocli_repro/cli_args.py:178-191, src/reprocli_repro/cli_args.py:56-60,
% src: appendix/a_harness.tex:22-25, appendix/a_harness.tex:104-107
The loop limits are inherited rather than set per run.
\texttt{--tool-rounds} carries the round cap and
\texttt{--budget-h100-hours} carries a flat compute ceiling. Leaving the
budget flag off is the configuration every sweep in this paper used, and each
paper's ceiling is then derived from the upper edge of its compute band.
Appendix~\ref{app:harness} gives both values and the metering rules they
enforce. Sampling fields are left unset, so the served model's own generation
configuration applies.

% src: src/reprocli_repro/cli_args.py:88-127, README.md:88-100,
% src: scripts/reproduce/repro_audit_one.sh:1-20
\begin{codebox}[lst:running-one]{bash}{Running one paper (sanitized)}
# One paper of the frozen benchmark, against an already-served model.
$ PYTHONPATH=src python3 -m reprocli_repro \
    --paper-id <arxiv-id> \
    --split eval \
    --run-id <run-id> \
    --runs-dir <runs-dir> \
    --vllm-server-url <endpoint-url> \
    --served-model-name <model-id>

# The dev split, for a smoke test that is not part of the benchmark.
$ PYTHONPATH=src python3 -m reprocli_repro \
    --paper-id 2505.18513 --split dev \
    --vllm-server-url <endpoint-url>

# Reproduce and grade the same bundle in one step.
$ scripts/reproduce/repro_audit_one.sh <arxiv-id> eval
\end{codebox}

% src: src/reprocli_repro/inputs.py:42-49, src/reprocli_repro/batch_runs.py:112-125,
% src: README.md:102-106
The run writes its bundle to
\texttt{<runs-dir>/<arxiv\_id>/<budget>h/<run\_id>/}, holding
\texttt{report.json}, \texttt{evidence/}, and \texttt{stats.json} alongside
the \texttt{workspace/} and the read-only \texttt{reference/}.
Appendix~\ref{app:report-format} describes what each part contains. The budget
directory is the only level between the paper and the run, so naming a paper
and a run identifier resolves to exactly one bundle. Naming only a paper does
not, because a paper attempted more than once has several.

% src: src/reprocli_repro/__main__.py:9-15, src/reprocli_repro/__main__.py:118-131,
% src: src/reprocli_vllm/vllm/endpoint.py:14-18
With no endpoint configured the same command becomes a dry run. It still loads
the lockfile row, prepares the episode's bundle, materializes the reference
copy, and prints the rendered task prompt, and it issues no requests and
allocates no GPU. This is how the inputs of Appendix~\ref{app:prompts} are
inspected offline. The auditor has no dry run and exits with an error instead,
because there is nothing to grade without a model.

\subsection{Running a tier}
\label{app:running-tier}

% src: scripts/reproduce/dsv4_flash/easy_dsv4_flash.sbatch:58-120,
% src: scripts/reproduce/dsv4_flash/easy_dsv4_flash.sbatch:280-300,357-360
One Slurm job runs one tier with one model. The job asks for a GPU share sized
to the model, one, two, or four GPUs of a node across the roster, serves the
model on it, waits for the endpoint file, and then walks the tier's
papers through the paired flow. Papers are enumerated from the lockfile by
their \texttt{tier} field and sorted by compute ceiling ascending, so cheap
papers land results early and an expensive paper meeting the job's wall clock
wastes less of the allocation. Each paper is dispatched as a background shell
function, with \texttt{MAX\_PARALLEL} papers in flight at once. Every paper in
flight adds one just-in-time GPU allocation of its own, so the parallelism
setting is bounded by queue pressure rather than by the served model.

% src: scripts/reproduce/dsv4_flash/easy_dsv4_flash.sbatch:424-525
Each paper's worker does four things in order: probe the served model and wait
rather than spend the paper if it is unreachable, reproduce with a pinned run
identifier, bind a grading root whose single entry is a symbolic link to that
exact bundle, and grade it. The grading root matters because the auditor's
contract is \texttt{<runs-dir>/<paper\_id>}, which on the sweep's run tree
mixes every past attempt at that paper. Each worker appends one row to a status file
recording the reproduction exit code, the audit exit code, the upload outcome,
and the run identifier. A paper that produced no bundle is recorded as a
harness casualty and is not counted as a model failure.

% src: scripts/reproduce/dsv4_flash/easy_dsv4_flash.sbatch:119,430-434
\texttt{RESUME=1} is what makes a killed sweep restartable. The worker skips a
paper whose verdict file is already on disk and non-empty, so resubmitting the
job after a wall-clock kill picks up only the tail. Nothing is regraded, and a
paper that died without writing a verdict is retried. Pointing the resubmitted
job's sweep directory at its predecessor's is the whole resume protocol.

% src: scripts/reproduce/dsv4_flash/easy_dsv4_flash.sbatch:58-120,280-360,424-525 (sanitized excerpt)
\begin{codebox}[lst:running-tier]{bash}{Sanitized excerpt of one tier's batch script}
#SBATCH -A <account>
#SBATCH -p <partition>
#SBATCH --nodes=1
#SBATCH --ntasks=1
#SBATCH --cpus-per-task=144
#SBATCH --gpus-per-node=2
#SBATCH --gpu-bind=none
#SBATCH --mem=220G
#SBATCH --time=48:00:00

SPLIT="${SPLIT:-eval}"                   # eval -> the frozen 100-paper split
TIER="${TIER:-Easy}"                     # only this tier's rows
TOOL_ROUNDS="${TOOL_ROUNDS:-300}"        # reproduction agent rounds
AUDIT_TOOL_ROUNDS="${AUDIT_TOOL_ROUNDS:-40}"   # served in-sweep auditor
MAX_PARALLEL="${MAX_PARALLEL:-10}"       # papers in flight; each adds one GPU allocation
RESUME="${RESUME:-1}"                    # 1 -> skip papers that already have a verdict

# --- Step 1: serve the model, wait for the published endpoint --------------
python -m reprocli_serve \
  --model "${MODEL}" --served-model-name "${SERVED_NAME}" \
  --port "${PORT}" --tensor-parallel-size "${TP}" \
  --endpoint-file "${ENDPOINT_FILE}" &
export REPROCLI_ENDPOINT_FILE="${ENDPOINT_FILE}"

# --- Step 2: enumerate the tier, smallest compute ceiling first ------------
mapfile -t PAPER_IDS < <(python - "${LOCKFILE}" "${SPLIT}" "${TIER}" <<'PY'
...
PY
)

# --- Step 3: reproduce + grade each paper, paired on one bundle ------------
run_one() {
  local pid="$1"
  local verdicts="${SWEEP_DIR}/audit_${pid}_extracted.jsonl"
  if [[ "${RESUME}" == "1" && -s "${verdicts}" ]]; then
    echo "=== ${pid} already graded -> skip (${verdicts})"
    return 0
  fi
  local rid="${OUTER_JOB}-${pid}-$(od -An -N3 -tx1 /dev/urandom | tr -d ' \n')"
  python -m reprocli_repro --paper-id "${pid}" --split "${SPLIT}" \
    --lockfile "${LOCKFILE}" --run-id "${rid}" \
    --tool-rounds "${TOOL_ROUNDS}" --served-model-name "${SERVED_NAME}" >"${log}" 2>&1
  local run_dir
  run_dir="$(find "${RUNS_DIR}/${pid}" -type d -name "${rid}" -print -quit)"
  ln -s "${run_dir}" "${grade_root}/${pid}"
  python3 src/run_arxiv_prompt_vllm.py --mode audit \
    --vllm-server-url "${BASE_URL}" --served-model-name "${AUDIT_MODEL}" \
    --runs-dir "${grade_root}" --claims "${CLAIMS}" \
    --paper-ids-file "${ids_file}" --tool-rounds "${AUDIT_TOOL_ROUNDS}" \
    --extracted-output "${verdicts}" >>"${log}" 2>&1
}

inflight=0
run_pids=()
for pid in "${PAPER_IDS[@]}"; do
  run_one "${pid}" &
  run_pids+=($!)
  inflight=$((inflight + 1))
  if (( inflight >= MAX_PARALLEL )); then
    wait -n
    inflight=$((inflight - 1))
  fi
done
wait "${run_pids[@]}"
\end{codebox}

% src: scripts/reproduce/dsv4_flash/easy_dsv4_flash.sbatch:58-120,280-360,424-525
The excerpt keeps the four structural steps and drops the model-specific serve
environment, the telemetry beacons, and the summary block.

\subsection{Grading a run}
\label{app:running-grade}

% src: src/reprocli_claude/__main__.py:1-25, src/reprocli_claude/__main__.py:57-104,
% src: src/reprocli_claude/agent.py:38-41, tasks/medium-tier-mechanism-synthesis.md:22
Every headline grade in this paper comes from one pinned grader.
\texttt{reprocli\_claude} runs the auditor of Appendix~\ref{app:rubric} over
the Anthropic Messages API with \texttt{--model claude-sonnet-5}. Everything
except the transport is shared with the served auditor: the same prompt
template, the same frozen rubric, the same pinned success bar from the
lockfile, the same run-directory tools, and the same finalizer that derives
the verdict and applies the deterministic caps. Its grader defaults are
\texttt{--effort high}, \texttt{--tool-rounds 25}, \texttt{--max-tokens
32000}, and \texttt{--workers 4}.

% src: src/reprocli_repro/batch_runs.py:1-20, src/reprocli_claude/__main__.py:106-140
Runs are selected by sweep or by paper. \texttt{--batch slurm-<jobid>}
resolves every run the batch script launched, takes each paper's newest
attempt, and binds it to the bundle its run identifier wrote.
Runs still in progress are dropped unless \texttt{--include-running} is
passed, because a half-written bundle grades the harness. \texttt{--paper-ids}
and \texttt{--paper-ids-file} name papers instead and fall back to each
paper's newest bundle that holds a run record. \texttt{--dry-run} prints the
resolved paper, run identifier, and bundle for each selection and stops.

% src: src/reprocli_claude/__main__.py:96-104, src/reprocli_claude/__main__.py:160-175,
% src: src/reprocli_claude/__main__.py:186-196, src/reprocli_claude/__main__.py:245-260
Each graded run produces two rows. \texttt{--output} receives the raw
response with its tool loop and token usage, and
\texttt{--extracted-output} receives the finalized verdict row with the score,
the verdict, the reproduced flag, the rationale, and any cheat flags. Each
grading pass carries an attempt tag of the grader model name joined to a UTC
timestamp, which keeps a re-grade of the same run distinct from the first
pass instead of colliding with rows already written. \texttt{--upload} writes
the verdict onto the run's row in the telemetry store, overwriting any prior
audit of that run, and it needs the store's URL and service key in the
environment. \texttt{--skip-audited} leaves already-graded runs alone, so a
repair pass only pays for what is missing.

% src: src/reprocli_claude/__main__.py:57-104, src/run_arxiv_prompt_vllm.py:1-60,
% src: scripts/reproduce/repro_audit_one.sh:83-100
\begin{codebox}[lst:running-grade]{bash}{Grading with the pinned grader (sanitized)}
# Re-grade a finished sweep with the pinned grader and file the verdicts.
$ PYTHONPATH=src python3 -m reprocli_claude \
    --batch slurm-<jobid> \
    --runs-dir <runs-dir> \
    --split eval \
    --model claude-sonnet-5 \
    --extracted-output <verdicts.jsonl> \
    --upload

# Resolve what would be graded, without calling the grader.
$ PYTHONPATH=src python3 -m reprocli_claude --batch slurm-<jobid> \
    --runs-dir <runs-dir> --dry-run
\end{codebox}

% src: src/run_arxiv_prompt_vllm.py:1-60, scripts/reproduce/repro_audit_one.sh:83-100, scripts/reproduce/dsv4_flash/easy_dsv4_flash.sbatch:482-500
The served auditor is the same grading procedure over a vLLM endpoint, run
with the tier script's \texttt{AUDIT\_TOOL\_ROUNDS} of 40 rounds where the
pinned grader has 25, invoked as
\texttt{python3 src/run\_arxiv\_prompt\_vllm.py --mode audit --claims
<claims-file> --runs-dir <runs-dir> --vllm-server-url <endpoint-url>}, and it
is what the tier batch scripts call inline while a sweep is still running.

\subsection{Reproducing the tables and listings}
\label{app:running-tables}

% src: .claude/skills/analyze-sweep/SKILL.md:52-60, .claude/skills/analyze-sweep/driver.py,
% src: src/reprocli_repro/batch_runs.py:1-20
Reconstructing a results table needs the per-sweep dumps and the evidence
bundles. Each sweep dumps to \texttt{runs.json}, one row per run carrying the
audit score, the verdict, the rationale, the cheat flags, the exit reason, the
spend against the budget, and the token counts. Alongside it,
\texttt{events/<run\_id>.jsonl} holds that run's ordered transcript and
\texttt{aggregates.json} the per-sweep totals. The per-run mechanism labels of
Appendix~\ref{app:taxonomy} live in \texttt{analyses.json}, one entry per run,
keyed by the same run identifiers. Together with the evidence bundles on disk,
those four files are enough to recompute every cell without a model.

% src: paper_latex/tables/gen_tool_schemas.py:1-18, paper_latex/tables/gen_audit_schema.py:1-7,
% src: paper_latex/tables/gen_audit_finalizer.py:1-11, paper_latex/tables/gen_classifier_schema.py:1-12,
% src: paper_latex/tables/gen_example_row.py:1-18, paper_latex/tables/gen_eval100.py:1-18
The appendix listings that quote machine-readable artifacts are generated
rather than transcribed, so a change in the harness is visible in the paper.
Each generator names its source paths at the top of the script, reads them
live from the repository, and writes into \texttt{paper\_latex/prompts/} or
\texttt{paper\_latex/tables/}. Table~\ref{tab:generators} maps each script to
what it regenerates. Two of them read the frozen split files rather than the
source tree, and \texttt{RECLAIM\_LOCKFILE\_DIR} points them at a different
copy.

% src: paper_latex/tables/gen_tool_schemas.py:1-18, paper_latex/tables/gen_audit_schema.py:1-7,
% src: paper_latex/tables/gen_audit_finalizer.py:1-11, paper_latex/tables/gen_classifier_schema.py:1-12,
% src: paper_latex/tables/gen_example_row.py:1-18, paper_latex/tables/gen_eval100.py:1-18
\begin{table}[h]
\caption{Generators for the appendix listings and tables.}
\label{tab:generators}
\begin{center}
\small
\begin{tabular}{@{}lp{8.3cm}@{}}
\toprule
Script under \texttt{paper\_latex/tables/} & What it regenerates \\
\midrule
\texttt{gen\_tool\_schemas.py} & The advertised \texttt{run\_gpu} and
\texttt{workspace\_bash} tool schemas and the report schema, read out of the
harness (Appendix~\ref{app:toolset}, Appendix~\ref{app:report-format}). \\
\texttt{gen\_audit\_schema.py} & The auditor's structured output schema
(Appendix~\ref{app:rubric}). \\
\texttt{gen\_audit\_finalizer.py} & The deterministic verdict finalizer,
excerpted verbatim from the auditor (Appendix~\ref{app:rubric}). \\
\texttt{gen\_classifier\_schema.py} & The artifact classifier's strict
response format for its final tools-off request. \\
\texttt{gen\_example\_row.py} & One frozen lockfile row, field for field, with
the prose fields cut at a word boundary
(Appendix~\ref{app:dataset-schema}). \\
\texttt{gen\_eval100.py} & The per-paper anchor tables for the frozen
100-paper split and the 14-paper dev split, plus the tier, band, and match-bar
counts the prose asserts (Appendix~\ref{app:dataset-schema}). \\
\bottomrule
\end{tabular}
\end{center}
\end{table}

% src: paper_latex/tables/gen_tool_schemas.py:1-18, paper_latex/tables/gen_eval100.py:1-18
Running a generator and rebuilding the paper is the check that a listing still
matches the code it was taken from.

\subsection{Tests}
\label{app:running-tests}

% src: README.md:158-166, verified by running the suite on 2026-08-30
The test suite exercises the harness offline. It needs no GPU, no network, and
no served model, and it covers the tool loop, the sandbox argument shapes, the
GPU session accounting, the report and verdict schemas, the lockfile loaders,
the serve command builder, and the grader's cost model. Running it is the
fastest check that a fresh checkout is wired correctly.

% src: verified by running `PYTHONPATH=src python3 -m pytest tests -q` on 2026-08-30
\begin{codebox}[lst:running-tests]{bash}{Running the suite}
$ PYTHONPATH=src python3 -m pytest tests -q
540 passed, 1 skipped
\end{codebox}

% src: verified by running `PYTHONPATH=src python3 -m pytest tests -q` on 2026-08-30
The 540 passing tests are the offline contract. Anything that needs a GPU, a
Slurm allocation, or a served model sits outside the suite and is exercised by
the commands above.
